\documentclass[10pt,a4paper]{amsart}
\usepackage[margin=19mm]{geometry}
\usepackage{amsmath,amssymb,mathtools}
\usepackage[scr=rsfs,cal=euler]{mathalfa}
\usepackage{xfrac}
\usepackage[unicode,hidelinks,bookmarks=false]{hyperref}
\newcommand{\ie}{i.e.\ }
\newcommand{\cf}{cf.\ }
\newcommand{\resp}{respectively\ }
\newcommand{\Subsection}{\subsection}
\providecommand{\nobreakdash}{\nobreak\hbox{-}}
\setlength{\emergencystretch}{2em}
\multlinegap0pt
\usepackage{amssymb}
\usepackage{dsfont}
\usepackage{float}
\usepackage{mathtools}
\usepackage{xcolor}
\newtheorem{lemma}{Lemma}
\newcommand{\E}{\mathbb{E}}
\title{Inverse Random Source and Cauchy Problems for Semi-Discrete Stochastic Parabolic Equations in Arbitrary Dimensions}
\author{Rodrigo Lecaros, Ariel A. P\'erez, Manuel F. Prado}
\date{Source: arXiv:2509.03760v2 (CC BY 4.0)}
\begin{document}
\maketitle

\noindent\emph{Source and licence.} Inverse Random Source and Cauchy Problems for Semi-Discrete Stochastic Parabolic Equations in Arbitrary Dimensions. Version arXiv:2509.03760v2 (CC BY 4.0), distributed by arXiv under the Creative Commons Attribution 4.0 International licence, http://creativecommons.org/licenses/by/4.0/, as recorded in the arXiv licence metadata for this exact version and shown on the abstract page https://arxiv.org/abs/2509.03760v2. Full licence notice: \texttt{licenses/arxiv-2509.03760-CC-BY-4.0.txt}.\par\smallskip

\noindent\textit{Editorial scope.} Counted unit: the article's lemma lemmaoftheu (source0.tex lines 544--560). The article's complete original proof of it is delivered with it (source0.tex lines 561--654, one \texttt{proof} environment of 94 lines). No mathematics is written, repaired or re-derived: the statement and the proof are the article's own lines, in the article's own notation, and no retained line is substituted or re-flowed.  Retained uncounted as the source-local context the proof needs: lines 94--101.  Nothing was omitted from any retained span, so the package carries no omission record.  No retained line contains a citation, so the wrapper carries no bibliography and no bibliography entry.  No retained line contains a figure environment, an image-inclusion command or drawing code, so no image is delivered and the package declares no asset dependency.  Editorial formatting only, in addition: an AMS \texttt{amsart} preamble that restates the article's own declarations and macros used by the retained lines, namely the environments \texttt{lemma} on separate counters, together with the standard packages those lines need.  The retained environments print in this document as Lemma 1 in order of appearance, whereas the article prints the same items with the numbering of its own full document.  The complete native source of the article as distributed in the arXiv e-print for this exact version is bundled unchanged as \texttt{sources/184437/source0.tex}.
\begin{equation}\label{normal:function}
    \forall x\in\partial_{i}\mathcal{M}, \quad
    \nu_{i}(x):=\begin{cases} 
     1 &  \mbox{if }x-\tfrac{h}{2}e_i\in \mathcal{M}_{i}^{\ast} \ \mbox{and } x+\tfrac{h}{2}e_i\notin\mathcal{M}_{i}^{\ast},\\ 
    -1 &  \mbox{if }x-\tfrac{h}{2}e_i\notin \mathcal{M}_{i}^{\ast} \ \mbox{and }x+\tfrac{h}{2}e_i\in\mathcal{M}_{i}^{\ast},\\ 
     0 &  \mbox{otherwise}.
    \end{cases}
\end{equation}

\begin{lemma}\label{lemmaoftheu}
Let $\xi \in H_{\mathcal{F}}^1(0,T; H_h^{1/2}(\partial \mathcal{M}) \cap H_h^1(\partial \mathcal{M}))$ with $\left.\xi\right|_{t=0} = 0$. Then, for a.e. $\omega \in \Omega$, the solution $u \in L^2(0,T; H_h^2(\mathcal{M})) \cap C(0,T; L_h^2(\mathcal{M}))$ of the system:
\begin{equation}\label{sys:homogeneuous}
    \begin{cases}
    du - \sum_{i=1}^{n} D_i(\gamma_i D_i u) \, dt = 0, & \text{in } Q, \\
    u = \xi,\quad\text{on } \partial Q,\quad \left.u\right|_{t=0} = 0,\quad  \text{in } \mathcal{M},
    \end{cases}
\end{equation}
satisfies the following stability estimate for any constant $\mathcal{C} > 0$:
\begin{equation}\label{ine:stability:H2}
    \|u\|_{L^2_{\mathcal{F}}(0,T; H_h^2(\mathcal{M}))} + \|u\|_{L^2_{\mathcal{F}}(\Omega; C([0,T]; L_h^2(\mathcal{M})))} \leq \mathcal{C} \|\xi\|_{H^1_{\mathcal{F}}(0,T; H_h^{1/2}(\partial \mathcal{M}) \cap H_h^1(\partial \mathcal{M}))}.
\end{equation}
Moreover, the following boundary estimate holds for $t_r^i(D_i u)$ on $\partial\mathcal{M}$:
\begin{equation}\label{ine:stability:boundary}
    \sum_{i=1}^{n} \mathbb{E} \int_{\partial Q} t_r^i(|D_i u|^2) \, dt \leq \mathcal{C} \|\xi\|_{L^2_{\mathcal{F}}(0,T; H_h^{1/2}(\partial \mathcal{M}) \cap H_h^1(\partial \mathcal{M}))}.
\end{equation}
\end{lemma}

\begin{proof}
Noting that for a.e. $\omega \in \Omega$ and $t \in [0,T]$, $\xi \in H_h^{1/2}(\partial \mathcal{M})$, there exists $v \in H_h^1(\mathcal{M})$ such that for a.e. $\omega \in \Omega$ and $t \in [0,T]$, $v = \xi$ on $\partial \mathcal{M}$ and
\begin{equation}\label{eq:tracexi}
\|v\|_{H_h^1(\mathcal{M})} =  \|\xi\|_{H_h^{1/2}(\partial \mathcal{M})}.
\end{equation}
Let $\tilde{u} = u - v$. Since $v|_{t=0}=0$ by \eqref{eq:tracexi} and $\xi|_{t=0}=0$, it follows that $\tilde{u}|_{t=0} = u|_{t=0} - v|_{t=0} = 0$. Moreover, on $\partial\mathcal{M}\times(0,T)$ we have $\tilde{u} = \xi - \xi = 0$. Therefore, $\tilde{u}$ satisfies
\begin{equation}\label{proof:sys:homogeneuous}
    \begin{cases}
    d\tilde{u} - \sum_{i=1}^{n} D_i(\gamma_i D_i \tilde{u}) \, dt = \left(\partial_t v - \sum_{i=1}^{n} D_i(\gamma_i D_i v)\right) \, dt, & \text{in } \mathcal{M} \times (0,T), \\
    \tilde{u} = 0, \quad\quad\text{on } \partial \mathcal{M} \times (0,T),\quad\quad
    \left.\tilde{u}\right|_{t=0} = 0,\quad \text{in } \mathcal{M}.
    \end{cases}
\end{equation}
The task is now to find an energy estimate for $u$. Let us apply the It\^{o}'s formula to $\displaystyle \frac{1}{2}\tilde{u}^2$ yields $\displaystyle \frac{1}{2}d(\tilde{u}^2)=\tilde{u}d\tilde{u}+\frac{1}{2}|d\tilde{u}|^2$. As $\tilde{u}$ satisfies \eqref{proof:sys:homogeneuous}, we have 
\[
\frac{1}{2}d(\tilde{u}^2)=\sum_{i=1}^{n}\tilde{u}D_i(\gamma_iD_i\tilde{u})\,dt+\tilde{u}\partial_{t}v\,dt-\sum_{i=1}^{n}\tilde{u}D_i(\gamma_iD_iv)\,dt
\]
By integration over $\mathcal{M}$ and integration by parts to the first and third terms, we get
\[
\frac{1}{2}\int_{\mathcal{M}}d(\tilde{u}^{2})=-\sum_{i=1}^{n}\int_{\mathcal{M}_i^{\ast}}\gamma_i|D_i(\tilde{u})|^2\,dt+\int_{\mathcal{M}}\tilde{u}\partial_{t}v\,dt+\sum_{i=1}^{n}\int_{\mathcal{M}_{i}^{*}}\gamma_iD_i(\tilde{u})\;D_i(v)\,dt,
\]
which is due to the fact that $\tilde{u}=0$ on $\partial\mathcal{M}$. From Young's inequality on the second and third terms in the above equation, it follows that
\[
\frac{1}{2}\int_{\mathcal{M}}d(\tilde{u}^2)\leq-\frac{1}{2}\sum_{i=1}^{n}\int_{\mathcal{M}_i^{*}}\gamma_i|D_i\tilde{u}|^2\,dt+\frac{1}{2}\int_{\mathcal{M}}|\tilde{u}|^2\,dt+\frac{1}{2}\int_{\mathcal{M}}|\partial_{t}v|^2\,dt+\frac{1}{2}\sum_{i=1}^{n}\int_{\mathcal{M}_i^{*}}\gamma_i|D_iv|^2\,dt
\]
Now, integrating on $[0,t]$ for $t\in [0,T]$, using \eqref{eq:tracexi} of the above inequality and from condition that $\left.\tilde{u}\right|_{t=0}=0$, we can assert that
\begin{equation}\label{eqE}
    \frac{1}{2} \, \int_{\mathcal{M}}|\tilde{u}|^2+ \frac{1}{2}\sum_{i=1}^{n}\int_{Q_i^{*}} \gamma_i|D_i\tilde{u}|^2\, dt \leq \mathcal{C}\int_{0}^{T}\left(\|\xi_t\|^2_{ H_h^{1/2}(\partial \mathcal{M})}+\|\xi\|^2_{H_h^{1/2}(\partial \mathcal{M})}\right)\, dt.
\end{equation}
Recalling $u = \tilde{u} + v$ and using the above inequality, we deduce that
\[
\begin{aligned}
\frac{1}{2}\int_{\mathcal{M}}|u|^2+\frac{1}{2}\sum_{i=1}^{n}\int_{Q_i^{*}}\gamma_i|D_iu|^2\,dt\leq \mathcal{C}\int_{0}^{T}\left(\|\xi_t\|^2_{ H_h^{1/2}(\partial \mathcal{M})}+\|\xi\|^2_{H_h^{1/2}(\partial \mathcal{M})}\right)\, dt+\frac{1}{2}\int_{\mathcal{M}}|v|^2.
\end{aligned}
\]
To estimate the term $\displaystyle \int_{\mathcal{M}}|v|^2$, we note that applying Young's inequality and \eqref{eq:tracexi} enables us to write
\[
\frac{1}{2}\int_{\mathcal{M}}\partial_{t}|v|^2\leq \frac{1}{2}\int_{\mathcal{M}}v\,\partial_{t}v\leq \frac{1}{2}\int_{\mathcal{M}}|v|^2+|v_{t}|^2\leq \frac{1}{2}\left(\|\xi\|^2_{H^{1/2}_{h}(\partial \mathcal{M})}+\|\xi_{t}\|_{H^{1/2}_{h}(\partial\mathcal{M})}\right)  
\]
Thus, 
\[
\int_{\mathcal{M}}|v|^2\leq \int_0^{T}\|\xi_{t}\|^2_{H_{h}^{1/2}(\partial\mathcal{M})}+\|\xi\|^2_{H_{h}^{1/2}(\partial \mathcal{M})}\,dt
\]
where we used the integration on $[0,T]$ and that $\left.v\right|_{t=0}=0$.
Combining these inequalities, we conclude that
\begin{equation}\label{eq:estimateenergy}
  \frac{1}{2}\int_{\mathcal{M}}|u|^2+\frac{1}{2}\sum_{i=1}^{n}\int_{Q_i^{*}}\gamma_i|D_iu|^2\,dt\leq \mathcal{C}\int_{0}^{T}\left(\|\xi_t\|^2_{ H_h^{1/2}(\partial \mathcal{M})}+\|\xi\|^2_{H_h^{1/2}(\partial \mathcal{M})}\right)\, dt.
\end{equation}
Therefore, taking the supremum over $[0,T]$ in the above inequality only considered the first terms, we conclude that
\begin{equation*}
\|u\|_{C(0,T; L_h^2(\mathcal{M}))} \leq \mathcal{C}\int_{0}^{T} \|\xi_{t}\|^2_{H^{1/2}_{h}(\partial \mathcal{M})}+\|\xi\|^2_{H_{h}^{1/2}(\partial\mathcal{M})} dt.
\end{equation*}
Furthermore, using \eqref{eq:estimateenergy}, the above inequality and integrating over $(0,T)$, we have
\begin{equation*}
\|u\|_{L_h^2(0,T; H^1_{h}(\mathcal{M}))} \leq\mathcal{C}\int_{0}^{T} \|\xi_{t}\|^2_{H^{1/2}_{h}(\partial \mathcal{M})}+\|\xi\|^2_{H_{h}^{1/2}(\partial\mathcal{M})} dt.
\end{equation*}
Similar arguments apply to the case $D_iu$ for each $i=1,\ldots,n$, yields
\[
\|u\|_{L_h^2(0,T; H_{h}^2(\mathcal{M}))} \leq\mathcal{C}\int_{0}^{T} \|\xi_{t}\|^2_{H^{1/2}_{h}(\partial \mathcal{M})\cap H^{1}(\partial\mathcal{M})}+\|\xi\|^2_{H_{h}^{1/2}(\partial\mathcal{M})\cap H^1_{h}(\partial\mathcal{M})} dt.
\]
Combining these inequalities and taking the expectation, we derive the first estimate \eqref{ine:stability:H2}.

%%%%%%%%%%%%%%%%%%%%%%%%%%%%%%%%%%%%%%%%%%%%%%%%%%%%%%%%%%%%%%%%%%%%%%%%%%%%%%%%%%%%%%%%%%%%%%%%%%%%%%%%%%%%%%%%%%%%%%%%%%%%%%%%%%%%%%%%%%%%%%%%%%%%%%%%
Our next goal is to determine the estimate for $\displaystyle \sum_{i=1}^{n}\E\int_{\partial Q}t_{r}^{i}(|D_iu|^2)\,dt$. In order to obtain this inequality, it is convenient to consider the following 
\begin{equation}\label{eq:identityD}
2x_{j}D_{j}^2(u)A_{j}D_ju=x_{j}D_{j}(|D_ju|^2)
\end{equation}
where $x_{j}=x\cdot e_{j}$.

Integrating the results over $\mathcal{M}\times(0,T)$ and applying an integration by parts it follows
\begin{align*}
    \frac{1}{2}\int_{\mathcal{M}}x_{j}D_{j}(|D_{j}u|^{2})=&-\frac{1}{2}\int_{\mathcal{M}^{\ast}_{j}}D_{j}(x_{j})|D_{j}u|^{2}+\frac{1}{2}\int_{\partial_{j}\mathcal{M}}x_{j}t_{r}^{j}(|D_{j}u|^{2})\nu_{j}.
\end{align*}
Thus, denoting $\Gamma_{j}:=\{ x\in\partial_{j}\mathcal{M}\mid x_{j}\nu_{j}= 1\}\subset\partial_j\mathcal{M}$ it follows that the above equals
\begin{align}\label{eq:boundary:1}
    \frac{1}{2}\int_{\mathcal{M}}x_{j}D_{j}(|D_{j}u|^{2})=-\frac{1}{2}\int_{\mathcal{M}_{j}^{\ast}}|D_{j}u|^2+\frac{1}{2}\int_{\Gamma_{j}}t_{r}^{j}(|D_{j}u|^{2}).
\end{align}
Hence, integrating over $\mathcal{M}\times (0,T)$ in \eqref{eq:identityD}, and then  using equation \eqref{eq:boundary:1} we obtain
\begin{equation}
\begin{aligned}
\int_{Q}\,x_{j}D_j^2(u)A_{j}D_{j}(u)\,\,dt=-\frac{1}{2}\int_{\mathcal{M}_{j}^{\ast}}|D_{j}u|^2+&\frac{1}{2}\int_{\Gamma_{j}\times(0,T)}t_{r}^{j}(|D_{j}u|^{2})\,dt
\end{aligned}
\end{equation}
Finally, using the Young inequality and taking expectation we conclude that there exists a constant $\mathcal{C}$ independent of $h$ such that
\begin{equation}\label{ine:boundary:H^{2}}
    \E\int_{\Gamma_{j}\times(0,T)}t_{r}^{j}(|D_{j}u|^{2})dt\leq \mathcal{C}\left\| u\right\|^{2}_{L^{2}_{\mathcal{F}}(0,T;H_h^{2}(\mathcal{M}))}.
\end{equation}
The inequality \eqref{ine:boundary:H^{2}} stands for a bound for the elements of the boundary $\Gamma_{j}$, in the direction $x_{j}$, such that its normal function is equal to one. Similarly, it is possible to obtain a bound for the normal derivative on the boundary nodes where the normal function defined in \eqref{normal:function} is minus one, that is, to adopt the previous methodology by using the term $(1-x_{j})A_{j}D_{j}$ instead of the used to obtain the equation \eqref{eq:identityD}. Moreover, we emphasize that we can obtain the inequality \eqref{ine:boundary:H^{2}} in any direction. In conclusion, the solution of the system \eqref{sys:homogeneuous} satisfies
\begin{equation}\label{ine:boundary:normal:derivative}
   \sum_{i=1}^{n}\E\int_{\partial Q}t_{r}^{i}(|D_{i}u|^{2})\,dt\leq 
   \mathcal{C} \left\| u\right\|^{2}_{L_{\mathcal{F}}^{2}(0,T;H_h^{2}(\mathcal{M}))}. 
\end{equation}
Overall, combining \eqref{ine:stability:H2} and \eqref{ine:boundary:normal:derivative} we obtain \eqref{ine:stability:boundary}.
\end{proof}

\end{document}
